\documentclass[11pt]{article}
\usepackage[T1]{fontenc}
\usepackage{lmodern,amsmath,amssymb,amsthm,booktabs,microtype}
\usepackage[margin=.9in,top=.75in,bottom=.8in]{geometry}
\usepackage[hidelinks]{hyperref}
\newtheorem{theorem}{Theorem}
\newcommand{\oeis}[1]{\href{https://oeis.org/A#1}{A#1}}
\title{Diagonals and populations of five von Neumann cellular automata}
\author{Blake Stone mathematics repository\\[-2pt]\small Proof developed with OpenAI Codex}
\date{}
\begin{document}
\maketitle
\begin{abstract}
Starting from one live cell on the square lattice, we determine the diagonals
of rules 62, 65 and 73 at every stage. Rules 65, 321 and 323 share an
alternating cross whose geometry gives exact population formulas.
The results account for eighteen OEIS sequences. All sixteen conjectured
generating functions in the selected entries hold. Ten stated homogeneous
recurrences require later starting indices.
\end{abstract}
\section{Model and diagonal words}
At time zero only the origin is on. Write $c,e,w,u,d\in\{0,1\}$ for a site's
value and the values at its four nearest lattice neighbors. Rule $q$
outputs bit $c+2(e+w+u+d)$ of $q$, with bit zero least significant.
The displayed stage is the window $[-n,n]^2$ at time $n$.
The center-to-corner word reads $(0,0),(1,1),\ldots,(n,n)$, in this order.
Exponents on digits mean repetition.
\begin{theorem}\label{words}
The center-to-corner words are as follows.
\[
\begin{array}{ccl}
\text{rule}&\text{stages}&\text{word}\\\hline
62& n\ge2&0\,1^{\lfloor n/2\rfloor}0^{\lceil n/2\rceil}\\
65& n\ge5,\ n\text{ even}&11\,0^{n-1}\\
65& n\ge5,\ n\text{ odd}&00\,1^{n-1}\\
73& n\ge4,\ n\text{ even}&0^{n+1}\\
73& n\ge4,\ n\text{ odd}&1^{n+1}
\end{array}
\]
Rule 62 starts with $1,10$. Rules 65 and 73 start with $1,01,100,0011$.
At time four the rule-65 word is $10000$.
\end{theorem}
The stage-one rule-62 diagonal is $10$ while the horizontal ray reads $11$.
The published Mathematica extraction in \oeis{285613} takes a horizontal slice.
We use the stated diagonal, which agrees with the listed terms.
The orientation and the choice between base two and decimal digit strings
both affect the resulting integer sequence.

Robert Price supplied the sequences. Colin Barker proposed the diagonal,
population and generating-function formulas in 2016--17. Lars Blomberg proposed
the dyadic population formula for rules 65 and 323. The last two sections
give the index mapping and the recurrence corrections.

\newpage
\section{Rule 62 near the diagonal and outward front}
By reflection put $x=|X|$, $y=|Y|$, $r=x+y$ and $h=n-r$.
The active weights for rule 62 are exactly $1,2,3,4,5$.
\begin{theorem}\label{strip}
At every stage $n\ge7$ the cell value in each region is given by the table.
\[
\begin{array}{c|l}
\text{region}&\text{condition for being on}\\\hline
h<0&\text{never}\\
0\le h\le1&\text{always}\\
2\le h\le3&x>0,\ y>0,\ r\text{ even}\\
h\ge4,\ |x-y|\le2&r>0,\ r\text{ even}
\end{array}
\]
In the remaining region the proof allows either bit.
\end{theorem}
\begin{proof}
Let $M(h,x,y)$ be the set of bits allowed by the table. For a target
$(x,y)$, write $x^- = |x-1|$ and similarly for $y^-$. Its four source neighbors
have coordinates $(x+1,y),(x^-,y),(x,y+1),(x,y^-)$.
Evaluate their masks and the center mask at old time $n$. Let $L,H$
be the sums of their respective minima and maxima with weights $1,2,2,2,2$.
An on target is forced by $1\le L$ and $H\le5$. An off target is forced
by $H=0$ or $L\ge6$. These checks cover all assignments of unspecified bits.

Exchange $x,y$ so $x\le y$. With $d=y-x$ set
\[
 X'=\min(x,7),\quad
 Y'=X'+\begin{cases}d&d\le4,\\6+(d\bmod 2)&d>4,\end{cases}
 \qquad h'=\max(-2,\min(h,7)).
\]
The center and each of its four neighbors preserve the tests $r>0$,
$x,y>0$, parity of $r$, and $|x-y|\le2$. For $d>4$ every neighbor
stays outside the strip and parity is retained. For smaller $d$ the exact
gap is preserved. Truncating the smaller coordinate at 7 preserves the
zero tests. Subtracting an equal amount from both coordinates preserves
parity. Each delay tested is $h+j$ with $j\in\{-1,0,1\}$, for which truncation
preserves every threshold.
Enumerate $-2\le h'\le7$, $0\le X'\le7$, $0\le Y'\le14$. Retain triples with
$X'+Y'\ge3$ or $h'+X'+Y'\ge7$. Every image at time $n\ge7$ satisfies
this guard. If its compressed radius is below 3 the coordinates are unchanged
and truncating $h$ at 7 preserves the guard.
For each retained triple test $M(h'+1,X',Y')$ against $L,H$ as above.
All required inequalities hold. The executable table is
\texttt{certificate} in \texttt{Checkerboard.lean}. The base at time seven
is obtained by seven updates for $0\le x,y\le7$. Outside this box finite
propagation gives zero. Induction proves the assertion.
\end{proof}
On the diagonal $r=2x$, so for $n\ge7$ the bit is on precisely when
$1\le x\le\lfloor n/2\rfloor$. Six direct updates finish the early stages.

\newpage
\section{Alternating crosses and the uniform rule}
\begin{theorem}\label{cross}
For $q\in\{65,321,323\}$ the even stage $n=2m\ge6$ has its live cells on
the coordinate axes out to distance $m$, together with the four cells
$(\pm1,\pm1)$. The odd stage $n=2m+1\ge5$ is on outside the defect set
\[
\begin{split}
 &\{X=0,\ |Y|\le m+1\}\ \cup\ \{Y=0,\ |X|\le m+1\}\\
 {}\cup{}&\{|X|\le1,\ |Y|\le m\}\ \cup\ \{|Y|\le1,\ |X|\le m\}.
\end{split}
\]
Rules 65 and 321 coincide at every stage. Rule 323 differs only at stage 1.
Rule 73 is uniformly off at stage 4 and then alternates uniformly.
\end{theorem}
\begin{proof}
Call the indicated configurations $E_m$ and $O_m$. The following finite
test proves one update sends $E_m$ to $O_m$ and one sends $O_m$ to $E_{m+1}$
for all $m\ge3$. A coordinate $z\ge0$ is used only through $z=0$, $z=1$,
$z\le m$, $z\le m+1$ and $z\le m+2$. We need those tests also for its
neighbors $z+1$ and $|z-1|$. Set $M=\min(m,7)$. For $m\le7$ replace
$z$ by $Z=\min(z,11)$. For $m>7$ use
\[
 Z=\begin{cases}
 z&z\le3,\\
 4&z>3,\ z+2<m,\\
 \min(11,7+z-m)&\text{otherwise}.
 \end{cases}
\]
In the last branch $7+z-m\ge5$. The five comparisons, including at both
neighbors, have the same values for $(z,m)$ and $(Z,M)$.
It suffices to test $3\le M\le7$, $0\le X',Y'\le11$,
for each rule and transition. Direct substitution gives the claimed targets
throughout. The proof \texttt{tables} in \texttt{Cross.lean} checks all $720$
parameter and coordinate triples per rule.
Stage five gives $O_2$ and stage six gives $E_3$. Outside the
time-$n$ box all cells equal the uniform background by finite propagation.
The three rules agree at stages 0 and 2--5. At stage one rule 323 also
retains the origin, while rules 65 and 321 agree. For rule 73 the entire
time-four box is zero, as is the exterior. Its bits 0 and 9 are respectively 1
and 0, giving the uniform alternation thereafter.
\end{proof}
Reading the shapes on $X=Y$ proves the remaining cases of
Theorem~\ref{words}. At even times the population is $1+4m+4=2n+5$.
At odd times the defect count is
\[
  1+4(m+1)+4(2m-1)=12m+1.
\]
The terms count the origin, axial defects, and off-axis defects. The observed
population is $(2n+1)^2-(12m+1)=4n^2-2n+6$.
For rule 65 the five initial populations are $1,4,5,36,9$. For rule 323 they
are $1,5,5,36,9$. The growing observation window matters on odd stages.

\begin{center}
\begin{tabular}{c@{\qquad\qquad}c}
\setlength{\unitlength}{4.5pt}\begin{picture}(9,9)
\put(4,1){\rule{3.8pt}{3.8pt}}
\put(4,2){\rule{3.8pt}{3.8pt}}
\put(3,3){\rule{3.8pt}{3.8pt}}
\put(4,3){\rule{3.8pt}{3.8pt}}
\put(5,3){\rule{3.8pt}{3.8pt}}
\put(1,4){\rule{3.8pt}{3.8pt}}
\put(2,4){\rule{3.8pt}{3.8pt}}
\put(3,4){\rule{3.8pt}{3.8pt}}
\put(4,4){\rule{3.8pt}{3.8pt}}
\put(5,4){\rule{3.8pt}{3.8pt}}
\put(6,4){\rule{3.8pt}{3.8pt}}
\put(7,4){\rule{3.8pt}{3.8pt}}
\put(3,5){\rule{3.8pt}{3.8pt}}
\put(4,5){\rule{3.8pt}{3.8pt}}
\put(5,5){\rule{3.8pt}{3.8pt}}
\put(4,6){\rule{3.8pt}{3.8pt}}
\put(4,7){\rule{3.8pt}{3.8pt}}
\end{picture}&
\setlength{\unitlength}{4.5pt}\begin{picture}(9,9)
\put(0,0){\rule{3.8pt}{3.8pt}}
\put(1,0){\rule{3.8pt}{3.8pt}}
\put(2,0){\rule{3.8pt}{3.8pt}}
\put(3,0){\rule{3.8pt}{3.8pt}}
\put(5,0){\rule{3.8pt}{3.8pt}}
\put(6,0){\rule{3.8pt}{3.8pt}}
\put(7,0){\rule{3.8pt}{3.8pt}}
\put(8,0){\rule{3.8pt}{3.8pt}}
\put(0,1){\rule{3.8pt}{3.8pt}}
\put(1,1){\rule{3.8pt}{3.8pt}}
\put(2,1){\rule{3.8pt}{3.8pt}}
\put(6,1){\rule{3.8pt}{3.8pt}}
\put(7,1){\rule{3.8pt}{3.8pt}}
\put(8,1){\rule{3.8pt}{3.8pt}}
\put(0,2){\rule{3.8pt}{3.8pt}}
\put(1,2){\rule{3.8pt}{3.8pt}}
\put(2,2){\rule{3.8pt}{3.8pt}}
\put(6,2){\rule{3.8pt}{3.8pt}}
\put(7,2){\rule{3.8pt}{3.8pt}}
\put(8,2){\rule{3.8pt}{3.8pt}}
\put(0,3){\rule{3.8pt}{3.8pt}}
\put(8,3){\rule{3.8pt}{3.8pt}}
\put(0,5){\rule{3.8pt}{3.8pt}}
\put(8,5){\rule{3.8pt}{3.8pt}}
\put(0,6){\rule{3.8pt}{3.8pt}}
\put(1,6){\rule{3.8pt}{3.8pt}}
\put(2,6){\rule{3.8pt}{3.8pt}}
\put(6,6){\rule{3.8pt}{3.8pt}}
\put(7,6){\rule{3.8pt}{3.8pt}}
\put(8,6){\rule{3.8pt}{3.8pt}}
\put(0,7){\rule{3.8pt}{3.8pt}}
\put(1,7){\rule{3.8pt}{3.8pt}}
\put(2,7){\rule{3.8pt}{3.8pt}}
\put(6,7){\rule{3.8pt}{3.8pt}}
\put(7,7){\rule{3.8pt}{3.8pt}}
\put(8,7){\rule{3.8pt}{3.8pt}}
\put(0,8){\rule{3.8pt}{3.8pt}}
\put(1,8){\rule{3.8pt}{3.8pt}}
\put(2,8){\rule{3.8pt}{3.8pt}}
\put(3,8){\rule{3.8pt}{3.8pt}}
\put(5,8){\rule{3.8pt}{3.8pt}}
\put(6,8){\rule{3.8pt}{3.8pt}}
\put(7,8){\rule{3.8pt}{3.8pt}}
\put(8,8){\rule{3.8pt}{3.8pt}}
\end{picture}\\[-1pt]
\small Stage 6 & \small Stage 7
\end{tabular}\\[2pt]
{\small Rule 65 on $[-4,4]^2$. Black squares mark live cells.}
\end{center}

\newpage
\section{Ten diagonal readouts}
For $b\ge2$ write $R_b(k)=(b^k-1)/(b-1)$, $u=\lfloor n/2\rfloor$,
$v=\lceil n/2\rceil$. The forward and reverse values of a word
$0^f1^w0^t$ are $b^tR_b(w)$ and $b^fR_b(w)$. Theorem~\ref{words} gives
\[
\begin{array}{ccll}
\text{rule}&\text{stages}&\text{center to corner}&\text{corner to center}\\\hline
62&n\ge2& b^v R_b(u)& bR_b(u)\\
65&n\ge5,\ n\text{ even}&(b+1)b^{n-1}& b+1\\
65&n\ge5,\ n\text{ odd}&R_b(n-1)&b^2R_b(n-1)\\
73&n\ge4,\ n\text{ even}&0&0\\
73&n\ge4,\ n\text{ odd}&R_b(n+1)&R_b(n+1).
\end{array}
\]
Let $G(z)=\sum_{n\ge0}a_nz^n$. The last column lists every failure among
indices where the source claims the homogeneous recurrence for $Q$.
\begin{center}\small
\begin{tabular}{llllrl}
\toprule
Rule & $b$ & Reverse / forward OEIS & Denominator $Q$ & $\deg(QG)$ & Failures\\
\midrule
62 & 10 & \oeis{285612} (reverse) & $(1-z)(1-bz^2)$ &4&4\\
62 & 10 & \oeis{285613} (forward) & $(1-bz)(1-bz^2)$ &4&4\\
65 & 10 & \oeis{285643} / \oeis{285644} & $(1-z^2)(1-b^2z^2)$ &8&5, 6, 8\\
65 & 2  & \oeis{285645} / \oeis{285646} & same &8&5, 6, 8\\
73 & 10 & \oeis{285647} / \oeis{285648} & same &7&none\\
73 & 2  & \oeis{285649} / \oeis{285650} & same &7&none\\
\bottomrule
\end{tabular}
\end{center}
For rules 65 and 73 each parity subsequence eventually has the form
$A+B(b^2)^k$ at $n=2k$ or $2k+1$, yielding the indicated denominator.
For rule 62, $a_n-ba_{n-2}$ is constant in the reverse direction
for $n\ge4$. In the forward direction it is a constant multiple of $b^{n-1}$
for $n\ge4$. These give the displayed denominators. The numerator is obtained by
\[
                  [z^n](QG)=\sum_{j=0}^{\min(n,\deg Q)}Q_j a_{n-j}.
\]
For example \oeis{285613} has $QG=1-100z^2+1000z^4$. At $n=4$ the
recurrence predicts $100$ whereas the diagonal $01100$ gives $1100$.
For \oeis{285646} the numerator is $1+z-z^2-2z^3+4z^5+32z^6-32z^8$.
In particular $n=4$ does satisfy that recurrence. The three failures come later.
The coefficient lists for all ten are in the verification record.
They agree with the published numerators.

\newpage
\section{Populations, partial sums and first differences}
Write $p_n$ for a population and define
$d_n=p_{2^n-1}$, $s_n=\sum_{i=0}^n p_i$, $e_n=p_{n+1}-p_n$.
Let $\delta=0$ for rule 65 and $\delta=1$ for rule 323. From
Theorem~\ref{cross} and the five seed values we obtain
\[
\begin{aligned}
 d_n&=4\cdot4^n-10\cdot2^n+12 &&(n\ge2),\\
 6(s_n-\delta)&=
 \begin{cases}4n^3+35n-66&n\text{ even},\\
               4n^3+12n^2+35n-69&n\text{ odd}\end{cases} &&(n\ge4),\\
 e_n&=\begin{cases}4n^2+4n+3&n\text{ even},\\-4n^2+4n+1&n\text{ odd}\end{cases}
 &&(n\ge5).
\end{aligned}
\]
The first identity follows by substitution, including the exceptional stage
$2^2-1=3$. The partial-sum formula follows by taking successive differences
and matching $s_4=55+\delta$. The last identity follows by subtraction.
These verify the parity versions and the combined forms in the entries.
\begin{center}\small
\begin{tabular}{llllrl}
\toprule
Readout & Rule 65 & Rule 323 & Denominator $Q$ & $\deg(QG)$ & Failures in source\\
\midrule
$p_n$ & \oeis{269782}&\oeis{270569}&$(1-z^2)^3$&10&10 (rule 65)\\
$d_n$ & \oeis{270084}&\oeis{271254}&$(1-z)(1-2z)(1-4z)$&4&none\\
$s_n$ & \oeis{270085}&\oeis{271255}&$(1-z)^4(1+z)^3$&10&10 (rule 65)\\
$e_n$ & \oeis{270086}&\oeis{271256}&$(1-z)^2(1+z)^3$&9&9 (both)\\
\bottomrule
\end{tabular}
\end{center}
For rule 65 the population numerator is
\[
  N=1+4z+2z^2+24z^3-3z^4+4z^6+4z^7-8z^8+4z^{10}.
\]
Both parity subsequences are polynomials of degree at most two past their
seeds, so the sixth-order denominator annihilates the tail. Multiplying by it
gives $N$ from the initial values. For rule 323 add $z(1-z^2)^3$ to $N$.
Partial sums divide the series by $1-z$. For first differences the series is
$((1-z)G-p_0)/z$. For the dyadic sequence sum the three geometric progressions
and correct terms zero and one. This recovers the six published generating functions for
these eight entries. A270569 and A271255 supply the other two sequences and have
no listed conjectured generating function. The last nonzero degree in each
row gives the sharp eventual cutoff.

\paragraph{Sources and verification.}
All eighteen OEIS entries are linked in the tables. For the local rule convention
see S. Wolfram, \emph{A New Kind of Science} (2002), p.~170.
The Lean project checks the local dynamics and the infinite-time inductive
invariants by kernel reduction and integer arithmetic. The independent Python
verifier evolves finite exceptions against an infinite background and checks
the generating numerators against those printed in the entries.
\end{document}
